\section{Phép thử điều kiện cần}
\label{sec:ch08-dieu-kien-can}

\index{sanity check}%
\index{ngẫu nhiên hóa trọng số}%
Ba mục vừa rồi đều đi tìm một con số để xếp hạng các phương pháp. Adebayo và
cộng sự đi hướng ngược lại: thay vì hỏi phương pháp nào tốt hơn, họ hỏi phương
pháp nào đáng bị loại ngay, và họ trả lời bằng hai phép thử mà một lời giải
thích trung thực bắt buộc phải qua~\autocite{msanity}.

Cả hai đều là \tn{phép thử ngẫu nhiên hóa}{randomization test}, khác nhau ở chỗ
cái gì bị ngẫu nhiên hóa. Phép thử thứ nhất ngẫu nhiên hóa tham số mô hình: chạy phương pháp trên mô hình
đã huấn luyện, rồi chạy lại đúng phương pháp ấy trên một mạng cùng kiến trúc
nhưng khởi tạo ngẫu nhiên và chưa huấn luyện, và so hai kết quả. Phép thử thứ
hai ngẫu nhiên hóa dữ liệu: so kết quả trên mô hình huấn luyện với dữ liệu có
nhãn đúng, với kết quả trên cùng kiến trúc ấy huấn luyện trên một bản sao của dữ
liệu đã hoán vị ngẫu nhiên toàn bộ nhãn~\autocite{msanity}.

Lập luận đứng trên một quan hệ đơn giản. Một lời giải thích được cho là mô tả
điều mô hình đã học; nếu nó cho ra gần như cùng một bản đồ khi mô hình chưa học gì,
hoặc khi mô hình đã học một tương ứng ngẫu nhiên, thì nó đang mô tả một thứ
khác, nhiều khả năng là cấu trúc của chính đầu vào.

Điều hai phép thử ấy \emph{không} nói mới là chỗ quyết định cách đọc chúng, và
các tác giả phát biểu nó chính xác: một phương pháp trượt các phép thử này thì không đủ
sức phục vụ những nhiệm vụ đòi lời giải thích trung thực với mô hình hoặc với
quá trình sinh dữ liệu~\autocite{msanity}. Câu đó nói về chiều trượt, và chỉ
chiều trượt. Nó không nói rằng phương pháp qua được là trung thực, cũng không
cho một thang để xếp hai phương pháp cùng qua.

Vì vậy hai phép thử này không phải chỉ số theo
định nghĩa~\ref{def:ch08-chi-so-trung-thuc}: chúng không trả về một số mà quy
ước đọc là mức độ trung thực, mà trả về một phán quyết loại bỏ. Sự phân biệt
giữa điều kiện cần và mức độ nghe như chuyện thuật ngữ, nhưng nó đổi hẳn cách
đọc kết quả. Một phương pháp qua cả hai phép thử vẫn có thể sai hoàn toàn về cơ
chế, và không có gì trong công trình này ngăn điều đó.

Đặt cạnh nhau, bốn mục vừa rồi cho một bức tranh có hình dạng rõ. Họ xóa dần cho
số nhưng số ấy lẫn với phản ứng của mô hình trước đầu vào lạ. Huấn luyện lại gỡ
được chỗ lẫn ấy nhưng đổi luôn đại lượng được đo. Họ tiên đề đòi một đẳng thức
chặt hơn nhưng vẫn kiểm nó qua chính phép nhiễu đang bị nghi ngờ. Còn phép thử
điều kiện cần thì không đưa ra thang đo nào. Không có ai trong bốn nhóm ấy đối
chiếu con số của mình với ground truth độc lập, vì không ai có sẵn ground truth
như thế.
